% =====================================================================
% CHAPTER 9: Conclusions, Limitations and Future Works
% Template source: thesis-or-project-template_1745300785.pdf, Chapter 9.
% Assembled per thesis/RESTRUCTURE_PLAN.md from old mainmatter/ch5.tex:
% the three sections moved nearly verbatim under the template headings.
% Chapter references updated to the 9-chapter layout (results Chapter 8,
% optimization Chapter 7); all numbers and labels unchanged.
% =====================================================================
\chapter{Conclusions, Limitations and Future Works}

\section{Conclusion}
This thesis set out to build an edge-deployable deep stereo network
meeting two objectives. The first, an efficient AI-refined disparity
pipeline for resource-limited platforms, is met by StereoLite: 2.96
million parameters, 0.78 pixel end-point error on the SceneFlow test set,
10.9 percent D1-all zero-shot on Middlebury 2014, and 49.8 milliseconds
at half precision on a laptop RTX 3050. The optimized eight-bit form
measured 36.3 milliseconds on the Jetson Orin Nano, 27.5 frames per
second, near the real-time target. The second, tolerance to imperfect
rectification, is met by the perturbation study in Section
\ref{sec:rectification}: accuracy holds across the sub-pixel to
one-pixel vertical misalignment a calibrated rig produces, with
end-point error rising only from 1.03 to 1.53 pixels from the local
correlation lookup in the refinement stage.

The contributions of this thesis are the following.
\begin{itemize}
\item StereoLite, a 2.96 million parameter edge-budget stereo network,
combines a truncated YOLO feature encoder, a left-image recurrent
context stream, plane-aware coarse-to-fine tile refinement, and a
fail-soft quarter-resolution geometry encoding volume.
\item A controlled training pipeline combining native-crop sampling,
augmentation, and full-dataset training gives the compact tile-hypothesis
model its cross-domain accuracy, reaching 10.9 percent zero-shot
Middlebury 2014 D1-all, within four points of the larger
distillation-trained LiteAnyStereo.
\item An equivalence-proven inference-optimization pass cuts latency by a
factor of 1.74 with no loss of accuracy and identifies the operators
that block eight-bit deployment.
\item The model reaches 0.78 pixel end-point error on the SceneFlow test
set with real-time inference on a laptop GPU and transfers zero-shot to
an unseen benchmark and a real camera without adaptation.
\end{itemize}

\FloatBarrier
\section{Limitations}
This work has several limitations. Results come from a single training
run, so run-to-run variance is not characterized. A gap remains in
cross-domain accuracy: at 10.9 percent zero-shot D1-all the model trails
the distillation-trained LiteAnyStereo (6.9 percent) and the iterative
IGEV-Stereo (5.0 percent). The disparity search spans 192 pixels at the 384 by
640 inference resolution, matching the evaluation protocol; supervision
is synthetic only; and the model is not yet on the KITTI leaderboard. Architecturally, recurrent reasoning
stops at one quarter resolution, limiting recovery of thin and distant
structures, and the network lacks left-right consistency and occlusion
reasoning. The rectification tolerance measured in Chapter 8 holds only
for small vertical offsets, not for large de-rectification.

\FloatBarrier
\section{Future Works}
Several directions follow from these limitations, ordered by expected
return. The most promising is a three-stage knowledge-distillation
pipeline following LiteAnyStereo: synthetic training with dense
supervision as in this thesis, then self-distillation under input
perturbation with a frozen copy of the model supplying feature targets,
then distillation from a frozen foundation teacher on unlabeled real
pairs, using the teacher pipeline built for this project. Chapter 8's
cross-domain evidence points to the training pipeline as the largest
remaining lever: failure concentrates in a scene type real-data
distillation addresses, and the current supervised-only recipe leaves
it unused. A second direction is deployment hardening: the eight-bit
TensorRT engine already runs on the Jetson Orin Nano at 36.3 milliseconds
(Chapter 8), so the next step is a full accuracy audit of the quantized
engine against the full-precision model and a port to other edge boards
in the same class. A third is completing the zero-shot evaluation quartet with
KITTI and ETH3D to characterize generalization beyond the indoor domain,
across driving and mixed domains; both datasets are staged in the
project's evaluation infrastructure. A fourth is a temporal
extension that propagates tile-plane state across frames, since a
small ego-motion warps a slanted-plane hypothesis predictably and would
amortize initialization cost over the video stream. Finally, the
reconstruction path in Chapter 8 invites a mapping extension that fuses
per-frame point clouds into a persistent scene model. A compact
tile-hypothesis network trained with the distillation pipeline above
would serve as an off-the-shelf edge stereo component, the natural next
step for this work.
